% Đề xuất đề tài MixNext (ticket #10, spec §14). Biên dịch: tectonic reports/proposal.tex  (hoặc xelatex trên Overleaf)
\documentclass[11pt,a4paper]{article}
\usepackage[a4paper,margin=2cm]{geometry}
\usepackage{amsmath}
\usepackage{fontspec}
\setmainfont{texgyretermes}[Extension=.otf, UprightFont=*-regular, BoldFont=*-bold,
  ItalicFont=*-italic, BoldItalicFont=*-bolditalic]  % theo tên file: chạy được cả Tectonic lẫn Overleaf, có đủ dấu tiếng Việt
\usepackage{microtype}
\usepackage{enumitem}
\usepackage{booktabs}
\usepackage{tabularx}
\usepackage[hidelinks]{hyperref}

\setlist{nosep,leftmargin=1.4em}
\setlength{\parindent}{0pt}
\setlength{\parskip}{0.45em}
\pagestyle{empty}

\newcommand{\muc}[1]{\par\vspace{0.5em}{\large\bfseries #1}\par\vspace{0.15em}}

\begin{document}

\begin{center}
  {\small ĐỀ XUẤT ĐỀ TÀI BÀI TẬP LỚN · MÔN IT4613 · 11/2026}\\[0.5em]
  {\LARGE\bfseries MixNext}\\[0.3em]
  {\large Gợi ý bài hát tiếp theo vừa hợp gu vừa chuyển bài mượt}\\[0.2em]
  {\itshape Transition-aware Playlist Continuation}
\end{center}

\muc{1. Bài toán thực tế}
Khi nghe nhạc liên tục, như tính năng tự phát tiếp của ứng dụng nghe nhạc, nhạc nền ở quán cà phê, buổi tiệc hay phòng tập,
hệ thống chọn bài tiếp theo chủ yếu theo gu. Bài kế tiếp có thể hợp gu nhưng chênh nhịp, lệch tông hoặc tụt năng lượng
đột ngột, làm mạch nghe bị đứt. Đây là bài toán \emph{automatic playlist continuation} (RecSys Challenge 2018) với thêm một
mục tiêu thứ hai: bài tiếp theo phải nối được với bài đang phát.

\begin{itemize}
  \item \textbf{Mục đích:} gợi ý bài tiếp theo vừa đúng gu của phiên nghe, vừa nối mượt với bài đang phát; đo xem muốn
        chuyển bài mượt hơn thì phải trả giá bao nhiêu độ chính xác.
  \item \textbf{Yêu cầu:} gợi ý ngay từ vài bài đầu của phiên; mỗi gợi ý có lý do; người dùng chỉnh được mức ưu tiên giữa
        “đúng gu” và “mượt”; ứng dụng demo chạy trên laptop, không cần mạng.
  \item \textbf{Kịch bản:} người dùng chọn vài bài đang nghe; hệ thống trả hàng đợi 10 bài, mỗi bài kèm lý do như
        “cùng tông 8A, BPM 124 lên 126”, và người dùng tự chọn bài phát tiếp từ hàng đợi. Khi mở nhạc cho buổi tiệc,
        người dùng kéo thanh trượt về phía “mượt” và hàng đợi thay đổi ngay.
\end{itemize}

\muc{2. Phương pháp dự kiến}
Hệ gợi ý lai hai tầng.
\begin{itemize}
  \item \textbf{Tầng 1 — điểm gu:} sinh 100 ứng viên theo gu của phiên bằng lọc cộng tác (item-kNN theo bài liền kề,
        ALS trên ma trận playlist × bài), item2vec và mô hình tuần tự SASRec, học từ dữ liệu playlist. So với hai baseline
        là Popularity và item-kNN.
  \item \textbf{Tầng 2 — điểm mix:} xếp lại ứng viên theo độ chuyển bài từ bài đang phát. Có hai cách tính:
        (a) luật của DJ: độ lệch BPM (kể cả nhịp gấp đôi/một nửa), vòng tông Camelot, chênh lệch năng lượng;
        (b) bộ mã hóa chuyển bài học từ các cặp bài liền nhau trong playlist bằng học tương phản (InfoNCE).
  \item \textbf{Kết hợp:} điểm cuối $= \alpha \cdot \text{điểm gu (chuẩn hóa)} + (1-\alpha) \cdot \text{điểm mix}$,
        $\alpha \in [0, 1]$; trọng số các thành phần mix và $\alpha$ chọn bằng grid search trên tập validation.
  \item \textbf{Đánh giá:} độ chính xác bằng HitRate@10, NDCG@10, MRR; độ mượt trên chuỗi 10 bài bằng độ lệch BPM,
        tỉ lệ cặp chuyển tông tương thích, độ lệch năng lượng và tỉ lệ cặp mix được. Kết quả chính là
        \textbf{đường đánh đổi} giữa độ chính xác và độ mượt khi thay đổi $\alpha$, cùng ablation từng thành phần điểm mix
        và so sánh luật với điểm mix học, lặp 3 lần chạy.
\end{itemize}

\muc{3. Dữ liệu}
Spotify Million Playlist Dataset (RecSys Challenge 2018; bản sao công khai trên Kaggle vì bản gốc đã bị gỡ), ghép với
đặc trưng âm thanh (tempo, tông, mode, năng lượng) từ một bộ dữ liệu công khai khoảng 2 triệu bài của MPD theo mã bài hát.
Mỗi playlist là một phiên nghe. Từ khoảng 100 nghìn playlist lấy mẫu, bỏ các bài không có đặc trưng hoặc xuất hiện dưới
5 lượt, bỏ các phiên quá ngắn, còn khoảng 92 nghìn phiên; chia train/validation/test 80/10/10 theo phiên. Khi đánh giá,
5 bài đầu của phiên là seed và bài thứ 6 là bài đích.

\muc{4. Đầu vào và đầu ra}
\begin{itemize}
  \item \textbf{Đầu vào:} danh sách một hoặc vài bài đang nghe (seed); tham số $\alpha$ từ 0 đến 1; số bài cần gợi ý $K$.
  \item \textbf{Đầu ra:} danh sách $K$ bài xếp hạng, mỗi bài có điểm gu, điểm mix và một dòng giải thích. Ở chế độ so sánh:
        chuỗi $K$ bài liên tiếp kèm biểu đồ BPM và năng lượng theo thứ tự bài, đặt cạnh chuỗi chỉ theo gu.
\end{itemize}

\muc{5. Thành viên}
\begin{tabularx}{\linewidth}{@{}p{0.3\linewidth}p{0.25\linewidth}X@{}}
  \toprule
  \textbf{Họ tên} & \textbf{Mã số sinh viên} & \textbf{Email} \\
  \midrule
  Nguyễn Hoàng Vũ  & 20235251 & \href{mailto:vu.nh235251@sis.hust.edu.vn}{vu.nh235251@sis.hust.edu.vn} \\
  Bùi Công Hoan    & 20235083 & \href{mailto:hoan.bc235083@sis.hust.edu.vn}{hoan.bc235083@sis.hust.edu.vn} \\
  Nguyễn Công Minh & 20235161 & \href{mailto:minh.nc235161@sis.hust.edu.vn}{minh.nc235161@sis.hust.edu.vn} \\
  \bottomrule
\end{tabularx}

\end{document}
